\begingroup
\setlength{\LTcapwidth}{\textwidth}
% Lebar kolom diseimbangkan: kolom Dampak dilebarkan agar isinya tidak terpecah 4-5 baris.
\begin{longtable}{p{0.19\textwidth}p{0.39\textwidth}p{0.30\textwidth}}
\caption{Ringkasan keterbatasan penelitian dan dampaknya terhadap interpretasi hasil.}
\label{tab:tab6.21_ringkasan_keterbatasan_dan_dampak} \\
\toprule
\textbf{Aspek} & \textbf{Keterbatasan} & \textbf{Dampak} \\
\midrule
\endfirsthead

\multicolumn{3}{@{}l}{\small\emph{Tabel \thetable{} (lanjutan)}} \\
\toprule
\textbf{Aspek} & \textbf{Keterbatasan} & \textbf{Dampak} \\
\midrule
\endhead

\bottomrule
\endlastfoot

Dataset & OhioT1DM tidak mewakili seluruh populasi; hipoglikemia relatif jarang & Generalisasi dan sensitivitas kelas berisiko terbatas \\

Generalisasi & Validasi lintas-pasien belum sama dengan validasi populasi luas & Belum dapat mengklaim performa universal \\

Uncertainty & Artefak kalibrasi belum tersedia untuk finger-stick 240 menit & Interval tidak tersedia seragam pada semua jalur \\

Retrieval & Manfaat conditioning terutama pada kasus divergen dan bergantung pada kondisi prediksi & Tidak merupakan peningkatan universal \\

Generation & Konfigurasi evaluasi lama berbeda dari konfigurasi production terbaru & Angka generation harus dibatasi pada protokolnya \\

Ahli/Usability & n=12 dan pengalaman CDSS/CGM terbatas & Evaluasi bersifat awal, bukan validasi klinis luas \\

Deployment & Backend lokal melalui ngrok & Kelayakan bergantung pada perangkat, jaringan, dan layanan LLM \\

Safety & Prediksi tidak menjamin sensitivitas tinggi untuk hipoglikemia & Sistem tetap doctor-mediated \\
\end{longtable}
\endgroup
